% main.tex
% gs -dNOPAUSE -dBATCH -sDEVICE=png16m -sOutputFile=out%03d.png -r1000x1000 main.pdf
% ffmpeg -framerate 30 -i out%03d.png out.mp4
% https://tex.stackexchange.com/a/748667
\unprotect
    %%% Defined similarly to \obeyMPboxdepth, \ignoreMPboxdepth,
    %%% \obeyMPboxorigin, and \normalMPboxdepth
    \permanent\protected\def\strutMPboxdepth{%
        \let\meta_relocate_box\meta_strut_box_depth%
    }

    \protected\def\meta_strut_box_depth{%
        \setbox\b_meta_graphic\vpack{%
            \box\b_meta_graphic%
            %%% Add some depth to the box without affecting its size
            \nohrule height-\openstrutdepth width 0pt depth \openstrutdepth\relax%
        }
    }

    %%% The pagecolumns output routine alredy adds a \vfil elsewhere, so the one
    %%% here isn't needed and disrupts the depth added by \strutMPboxdepth.
    \let\page_otr_fill_page_unchecked\relax
\protect

%%% Comment the line below to confirm that an overfull vbox is generated;
%%% uncomment it to confirm that the overfull vbox message goes away.
\strutMPboxdepth
\setupbodyfont[pagella]


\definepapersize
  [youtubeHD]
  [width=1920bp,height=1080bp]
\setuppapersize[youtubeHD][youtubeHD]
\setuplayout[
    location=middle,
    width=middle,
    height=middle,
    topspace=0pt,
    backspace=0pt,
    header=0pt,
    footer=0pt
]
\defineframed[myframed][
    width=800pt,
    align=normal,
    offset=1cm,
    frame=off
]
\setupbodyfont[modern,24pt]
\ctxloadluafile{parametric.lua}
\processMPfigurefile{register_mp_cmd.mp}

\startMPinclusions
def render_theorem_box =
numeric w, h;
w := CurrentWidth;
h := CurrentHeight;

path p;
p := (-w/2, -h/2) -- (w/2, -h/2) -- (w/2, h/2) -- (-w/2, h/2) -- cycle;
setbounds currentpicture to p;

picture theorem;
theorem := textext("\getbuffer[BufA]");

pair center_anchor, topleft_anchor;
center_anchor := center theorem;
topleft_anchor := ulcorner theorem;

enddef;


def jasper_thin =
    withpen pencircle scaled .1cm
enddef ;

def jasper_thick =
    withpen pencircle scaled .2cm
enddef ;

def jasper_lightgray =
    withcolor (0.6)
enddef ;

def jasper_lightgraypenthin =
    jasper_lightgray withpen pencircle scaled .1cm
enddef ;

def jasper_lightgraypenthick =
    jasper_lightgray withpen pencircle scaled .2cm
enddef ;


def jasper_darkgray =
    withcolor (0.4)
enddef ;

def jasper_darkgraypenthin =
    jasper_darkgray withpen pencircle scaled .15cm
enddef ;

def jasper_darkgraypenthick =
    jasper_darkgray withpen pencircle scaled .3cm
enddef ;
\stopMPinclusions


\starttext

    \startbuffer[BufA]
        \startframed[myframed]
            {\bf Theorem 2.2.2} A circle cannot have more than one 
            centre.\par \color[white]{{\it Proof.} Let us suggest that the conclusion 
            of the theorem is false; in particular that a circle has 
            two or more centres. Let us denote just two of these centres
            \im{O} and \im{P}.\par According to AXIOM 1, there is 
            exactly one line that passes through points \im{O} and 
            \im{P}. Let \im{M} and \im{N} be the points of intersection 
            of this line with the circle.\par Since we have suggested
            that \im{O} is a centre of the circle, the segments 
            \im{NO} and \im{OM} are congruent as two radii of the 
            circle centred at \im{O}: \im{NO=OM}. Similarly,
            \im{NP=PM} as two radii of the circle centred at 
            \im{P}.\par Point \im{O} lies in the interior of 
            \im{NP}, hence \im{NO<NP}; similarly, \im{PM<OM}.
            Then we can write: \dm{NO<NP=PM<OM=NO}, which means 
            that \im{NO<NO}. This result contradicts to the axiom 
            of congruence of segments which asserts that each 
            segments is congruent to itself. Thus the suggestion 
            of the existence of two (\im{P} and \im{O}) or more 
            distinct centres is false; therefore a circle cannot 
            have more than one centre, {\bf Q.E.D.}}
        \stopframed
    \stopbuffer
    % initial theorem
    \dorecurse{7*30}{%
        \startMPpage
            render_theorem_box ; 
            numeric t;
            t := 0;
            pair anchor;
            anchor := (1 - t) * center_anchor + t * topleft_anchor + t * (-2cm, -9cm);
            label(theorem, (0,3cm) - anchor);
        \stopMPpage%
    }
    % Translate the theorem to the top right corner.
    \dorecurse{2*30+1}{%
        \startMPpage
            render_theorem_box ;
            numeric t;
            t := (#1-1)/60;
            pair anchor;
            anchor := (1 - t) * center_anchor + t * topleft_anchor + t * (-2cm, -9cm);
            label(theorem, (0,3cm) - anchor);
        \stopMPpage%
    }
    % Pause once the theorem is in place.
    \dorecurse{2*30}{%
        \startMPpage
            render_theorem_box ;
            numeric t;
            t := 1; 
            pair anchor;
            anchor := (1 - t) * center_anchor + t * topleft_anchor + t * (-2cm, -9cm);
            label(theorem, (0,3cm) - anchor);
        \stopMPpage%
    }
    % Pause once the theorem is in place.
    \dorecurse{2*30}{%
        \startMPpage
            render_theorem_box ;
            numeric t;
            t := 1; 
            pair anchor;
            anchor := (1 - t) * center_anchor + t * topleft_anchor + t * (-2cm, -9cm);
            label(theorem, (0,3cm) - anchor);
        \stopMPpage%
    }
    \startbuffer[BufA]
        \startframed[myframed]
            {\bf Theorem 2.2.2} A circle cannot have more than one 
            centre.\par {\it Proof.} \color[white]{Let us suggest that the conclusion 
            of the theorem is false; in particular that a circle has 
            two or more centres. Let us denote just two of these centres
            \im{O} and \im{P}.\par According to AXIOM 1, there is 
            exactly one line that passes through points \im{O} and 
            \im{P}. Let \im{M} and \im{N} be the points of intersection 
            of this line with the circle.\par Since we have suggested
            that \im{O} is a centre of the circle, the segments 
            \im{NO} and \im{OM} are congruent as two radii of the 
            circle centred at \im{O}: \im{NO=OM}. Similarly,
            \im{NP=PM} as two radii of the circle centred at 
            \im{P}.\par Point \im{O} lies in the interior of 
            \im{NP}, hence \im{NO<NP}; similarly, \im{PM<OM}.
            Then we can write: \dm{NO<NP=PM<OM=NO}, which means 
            that \im{NO<NO}. This result contradicts to the axiom 
            of congruence of segments which asserts that each 
            segments is congruent to itself. Thus the suggestion 
            of the existence of two (\im{P} and \im{O}) or more 
            distinct centres is false; therefore a circle cannot 
            have more than one centre, {\bf Q.E.D.}}
        \stopframed
    \stopbuffer
    % Proof.
    \dorecurse{3*30}{%
        \startMPpage
            render_theorem_box ;
            numeric t;
            t := 1; 
            pair anchor;
            anchor := (1 - t) * center_anchor + t * topleft_anchor + t * (-2cm, -9cm);
            label(theorem, (0,3cm) - anchor);
        \stopMPpage%
    }
    \startbuffer[BufA]
        \startframed[myframed]
            {\bf Theorem 2.2.2} A circle cannot have more than one 
            centre.\par {\it Proof.} Let us suggest that the conclusion 
            of the theorem is false; in particular that a circle has 
            two or more centres. \color[white]{Let us denote just two of these centres
            \im{O} and \im{P}.\par According to AXIOM 1, there is 
            exactly one line that passes through points \im{O} and 
            \im{P}. Let \im{M} and \im{N} be the points of intersection 
            of this line with the circle.\par Since we have suggested
            that \im{O} is a centre of the circle, the segments 
            \im{NO} and \im{OM} are congruent as two radii of the 
            circle centred at \im{O}: \im{NO=OM}. Similarly,
            \im{NP=PM} as two radii of the circle centred at 
            \im{P}.\par Point \im{O} lies in the interior of 
            \im{NP}, hence \im{NO<NP}; similarly, \im{PM<OM}.
            Then we can write: \dm{NO<NP=PM<OM=NO}, which means 
            that \im{NO<NO}. This result contradicts to the axiom 
            of congruence of segments which asserts that each 
            segments is congruent to itself. Thus the suggestion 
            of the existence of two (\im{P} and \im{O}) or more 
            distinct centres is false; therefore a circle cannot 
            have more than one centre, {\bf Q.E.D.}}
        \stopframed
    \stopbuffer
    % Let us suggest that the conclusion 
    % of the theorem is false; in particular that a circle has 
    % two or more centres.
    \dorecurse{8*30}{%
        \startMPpage
            render_theorem_box ;
            numeric t;
            t := 1; 
            pair anchor;
            anchor := (1 - t) * center_anchor + t * topleft_anchor + t * (-2cm, -9cm);
            label(theorem, (0,3cm) - anchor);
        \stopMPpage%
    }%
    \startMPinclusions
        def cx = -15 enddef ;
        def cy = 5 enddef ;
        def R = 7 enddef ;
        def circle_centre =
            jasper_point[% origin
                fx = "cx"
                ,fy = "cy"
            ] ;
            jasper_append_label[% origin label
                px="cx+cos(5*pi/4)"
                ,py="cy+sin(5*pi/4)"
                ,text="\im{O}"
            ] ;
        enddef ;
    \stopMPinclusions%
    % draw circle's centre
    \dorecurse{3*30}{%
        \startMPpage
            render_theorem_box ;
            numeric t;
            t := 1;
            pair anchor;
            anchor := (1 - t) * center_anchor + t * topleft_anchor + t * (-2cm, -9cm);
            label(theorem, (0,3cm) - anchor);
            circle_centre ;
            jasper_render_segments ;
        \stopMPpage%
    }
    % extend radius from circle's center
    \dorecurse{2*30+1}{%
        \startMPpage
            render_theorem_box ;
            numeric t;
            t := 1;
            pair anchor;
            anchor := (1 - t) * center_anchor + t * topleft_anchor + t * (-2cm, -9cm);
            label(theorem, (0,3cm) - anchor);
            circle_centre ;
            jasper_curve[% radius
                ustart="0"
                ,ustop="1"
                ,usamples="200"
                ,fx="cx+R*u*(#1-1)/60"
                ,fy = cy
                ,drawcommands="jasper_thin"
            ] ;
            jasper_point[% terminal
                fx = "cx + R*(#1-1)/60"
                ,fy = "cy"
                ,fz="1"
                ,drawcommands="withcolor black"
            ] ;
            jasper_render_segments ;
        \stopMPpage%
    }
    % pause after extending radius
    \dorecurse{2*30+1}{%
        \startMPpage
            render_theorem_box ;
            numeric t;
            t := 1;
            pair anchor;
            anchor := (1 - t) * center_anchor + t * topleft_anchor + t * (-2cm, -9cm);
            label(theorem, (0,3cm) - anchor);
            circle_centre ;
            jasper_curve[% radius
                ustart="0"
                ,ustop="1"
                ,usamples="200"
                ,fx="cx+R*u"
                ,fy = cy
                ,drawcommands="jasper_thin"
            ] ;
            jasper_point[% terminal
                fx = "cx + R"
                ,fy = "cy"
                ,fz="1"
                ,drawcommands="withcolor black"
            ] ;
            jasper_render_segments ;
        \stopMPpage%
    }
    % draw the circle
    \dorecurse{4*30+1}{%
        \startMPpage
            render_theorem_box ;
            numeric t;
            t := 1;
            pair anchor;
            anchor := (1 - t) * center_anchor + t * topleft_anchor + t * (-2cm, -9cm);
            label(theorem, (0,3cm) - anchor);
            circle_centre ;
            jasper_curve[% radius
                ustart="0",
                ustop="1",
                usamples="20",
                fx="cx + R*u*cos(tau*(#1-1)/120)",
                fy="cy + R*u*sin(tau*(#1-1)/120)"
                ,drawcommands="jasper_thin"
            ] ;
            jasper_point[% terminal
                fx = "cx + R*cos(tau*(#1-1)/120)"
                ,fy = "cy + R*sin(tau*(#1-1)/120)"
                ,fz = "1"
            ] ;
            jasper_curve[% arc
                ustart="0",
                ustop="1",
                usamples="200",
                fx="cx + R*cos(u*tau*(#1-1)/120)",
                fy="cy + R*sin(u*tau*(#1-1)/120)"
                ,drawcommands="jasper_thin"
            ] ;
            jasper_render_segments ;
        \stopMPpage%
    }%
    \startMPinclusions
        def full_circle =
            circle_centre ;
            jasper_curve[% circle
                ustart="0",
                ustop="1",
                usamples="200",
                fx="cx + R*cos(u*tau)",
                fy="cy + R*sin(u*tau)"
                ,drawcommands="jasper_thin"
            ] ;
        enddef ;
    \stopMPinclusions%
    % pause after drawing circle
    \dorecurse{4*30}{%
        \startMPpage
            render_theorem_box ;
            numeric t;
            t := 1;
            pair anchor;
            anchor := (1 - t) * center_anchor + t * topleft_anchor + t * (-2cm, -9cm);
            label(theorem, (0,3cm) - anchor);
            full_circle ;
            jasper_curve[% radius
                ustart="0",
                ustop="1",
                usamples="20",
                fx="cx + R*u*cos(tau)",
                fy="cy + R*u*sin(tau)"
                ,drawcommands="jasper_thin"
            ] ;
            jasper_point[% terminal of radius
                fx = "cx + R*cos(tau)"
                ,fy = "cy + R*sin(tau)"
                ,fz = "1"
                ,drawcommands="withcolor black"
            ] ;
            jasper_render_segments ;
        \stopMPpage%
    }
    % removed constructions, pause
    \dorecurse{4*30}{%
        \startMPpage
            render_theorem_box ;
            numeric t;
            t := 1;
            pair anchor;
            anchor := (1 - t) * center_anchor + t * topleft_anchor + t * (-2cm, -9cm);
            label(theorem, (0,3cm) - anchor);
            full_circle ;
            jasper_render_segments ;
        \stopMPpage%
    }
    \startbuffer[BufA]
        \startframed[myframed]
            {\bf Theorem 2.2.2} A circle cannot have more than one 
            centre.\par {\it Proof.} Let us suggest that the conclusion 
            of the theorem is false; in particular that a circle has 
            two or more centres. Let us denote just two of these centres
            \im{O} and \im{P}.\par \color[white]{According to AXIOM 1, there is 
            exactly one line that passes through points \im{O} and 
            \im{P}. Let \im{M} and \im{N} be the points of intersection 
            of this line with the circle.\par Since we have suggested
            that \im{O} is a centre of the circle, the segments 
            \im{NO} and \im{OM} are congruent as two radii of the 
            circle centred at \im{O}: \im{NO=OM}. Similarly,
            \im{NP=PM} as two radii of the circle centred at 
            \im{P}.\par Point \im{O} lies in the interior of 
            \im{NP}, hence \im{NO<NP}; similarly, \im{PM<OM}.
            Then we can write: \dm{NO<NP=PM<OM=NO}, which means 
            that \im{NO<NO}. This result contradicts to the axiom 
            of congruence of segments which asserts that each 
            segments is congruent to itself. Thus the suggestion 
            of the existence of two (\im{P} and \im{O}) or more 
            distinct centres is false; therefore a circle cannot 
            have more than one centre, {\bf Q.E.D.}}
        \stopframed
    \stopbuffer
    % Let us denote just two of these centres
    % \im{O} and \im{P}.
    \dorecurse{7*30}{%
        \startMPpage
            render_theorem_box ;
            numeric t;
            t := 1;
            pair anchor;
            anchor := (1 - t) * center_anchor + t * topleft_anchor + t * (-2cm, -9cm);
            label(theorem, (0,3cm) - anchor);
            full_circle ;
            jasper_render_segments ;
        \stopMPpage%
    }
    % move P into position
    \dorecurse{2*30}{%
        \startMPpage
            render_theorem_box ;
            numeric t;
            t := 1;
            pair anchor;
            anchor := (1 - t) * center_anchor + t * topleft_anchor + t * (-2cm, -9cm);
            label(theorem, (0,3cm) - anchor);
            full_circle ;
            jasper_point[
                fx="cx+2*#1/60"
                ,fy="cy"
            ] ;
            jasper_append_label[
                px="cx + 2*#1/60 + cos(5*pi/4)"
                ,py="cy + sin(5*pi/4)"
                ,text="\im{P}"
            ] ;
            jasper_render_segments ;
        \stopMPpage%
    }%
    \startMPinclusions
        def circle_OP =
            full_circle ;
            jasper_point[
                fx="cx+2"
                ,fy="cy"
            ] ;
            jasper_append_label[
                px="cx + 2 + cos(5*pi/4)"
                ,py="cy + sin(5*pi/4)"
                ,text="\im{P}"
            ] ;
        enddef ;
    \stopMPinclusions%
    % pause after moving P into position
    \dorecurse{4*30}{%
        \startMPpage
            render_theorem_box ;
            numeric t;
            t := 1;
            pair anchor;
            anchor := (1 - t) * center_anchor + t * topleft_anchor + t * (-2cm, -9cm);
            label(theorem, (0,3cm) - anchor);
            circle_OP ;
            jasper_render_segments ;
        \stopMPpage%
    }
    \startbuffer[BufA]
        \startframed[myframed]
            {\bf Theorem 2.2.2} A circle cannot have more than one 
            centre.\par {\it Proof.} Let us suggest that the conclusion 
            of the theorem is false; in particular that a circle has 
            two or more centres. Let us denote just two of these centres
            \im{O} and \im{P}.\par According to AXIOM 1, there is 
            exactly one line that passes through points \im{O} and 
            \im{P}. \color[white]{Let \im{M} and \im{N} be the points of intersection 
            of this line with the circle.\par Since we have suggested
            that \im{O} is a centre of the circle, the segments 
            \im{NO} and \im{OM} are congruent as two radii of the 
            circle centred at \im{O}: \im{NO=OM}. Similarly,
            \im{NP=PM} as two radii of the circle centred at 
            \im{P}.\par Point \im{O} lies in the interior of 
            \im{NP}, hence \im{NO<NP}; similarly, \im{PM<OM}.
            Then we can write: \dm{NO<NP=PM<OM=NO}, which means 
            that \im{NO<NO}. This result contradicts to the axiom 
            of congruence of segments which asserts that each 
            segments is congruent to itself. Thus the suggestion 
            of the existence of two (\im{P} and \im{O}) or more 
            distinct centres is false; therefore a circle cannot 
            have more than one centre, {\bf Q.E.D.}}
        \stopframed
    \stopbuffer
    % According to AXIOM 1, there is 
    % exactly one line that passes through points \im{O} and 
    % \im{P}.
    \dorecurse{8*30}{%
        \startMPpage
            render_theorem_box ;
            numeric t;
            t := 1;
            pair anchor;
            anchor := (1 - t) * center_anchor + t * topleft_anchor + t * (-2cm, -9cm);
            label(theorem, (0,3cm) - anchor);
            circle_OP ;
            jasper_render_segments ;
        \stopMPpage%
    }
    \startMPinclusions
        def line_circle =
            jasper_curve[
                ustart="-R-3"
                ,ustop="R+3"
                ,usamples="20"
                ,fx="cx+u"
                ,fy="cy"
                ,fz="-1"
                ,drawcommands="jasper_thin"
            ] ;
            circle_OP ;
        enddef ;
    \stopMPinclusions
    % draw the line
    \dorecurse{6*30}{%
        \startMPpage
            render_theorem_box ;
            numeric t;
            t := 1;
            pair anchor;
            anchor := (1 - t) * center_anchor + t * topleft_anchor + t * (-2cm, -9cm);
            label(theorem, (0,3cm) - anchor);
            line_circle ;
            jasper_render_segments ;
        \stopMPpage%
    }
    \startbuffer[BufA]
        \startframed[myframed]
            {\bf Theorem 2.2.2} A circle cannot have more than one 
            centre.\par {\it Proof.} Let us suggest that the conclusion 
            of the theorem is false; in particular that a circle has 
            two or more centres. Let us denote just two of these centres
            \im{O} and \im{P}.\par According to AXIOM 1, there is 
            exactly one line that passes through points \im{O} and 
            \im{P}. Let \im{M} and \im{N} be the points of intersection 
            of this line with the circle.\par \color[white]{Since we have suggested
            that \im{O} is a centre of the circle, the segments 
            \im{NO} and \im{OM} are congruent as two radii of the 
            circle centred at \im{O}: \im{NO=OM}. Similarly,
            \im{NP=PM} as two radii of the circle centred at 
            \im{P}.\par Point \im{O} lies in the interior of 
            \im{NP}, hence \im{NO<NP}; similarly, \im{PM<OM}.
            Then we can write: \dm{NO<NP=PM<OM=NO}, which means 
            that \im{NO<NO}. This result contradicts to the axiom 
            of congruence of segments which asserts that each 
            segments is congruent to itself. Thus the suggestion 
            of the existence of two (\im{P} and \im{O}) or more 
            distinct centres is false; therefore a circle cannot 
            have more than one centre, {\bf Q.E.D.}}
        \stopframed
    \stopbuffer
    % Let \im{M} and \im{N} be the points of intersection 
    % of this line with the circle.
    \dorecurse{10*30}{%
        \startMPpage
            render_theorem_box ;
            numeric t;
            t := 1;
            pair anchor;
            anchor := (1 - t) * center_anchor + t * topleft_anchor + t * (-2cm, -9cm);
            label(theorem, (0,3cm) - anchor);
            line_circle ;
            jasper_render_segments ;
        \stopMPpage%
    }%
    \startMPinclusions
        def circle_NM =
            line_circle ;
            jasper_point[
                fx="cx-R"
                ,fy="cy"
                ,fz="2"
            ] ;
            jasper_point[
                fx="cx+R"
                ,fy="cy"
                ,fz="2"
            ] ;
            jasper_append_label[
                px="cx-R+cos(5*pi/4)"
                ,py="cy+sin(5*pi/4)"
                ,text="\im{N}"
            ] ;
            jasper_append_label[
                px="cx+R+cos(5*pi/4)"
                ,py="cy+sin(5*pi/4)"
                ,text="\im{M}"
            ] ;
        enddef ;
    \stopMPinclusions%
    % add labels
    \dorecurse{7*30}{%
        \startMPpage
            render_theorem_box ;
            numeric t;
            t := 1;
            pair anchor;
            anchor := (1 - t) * center_anchor + t * topleft_anchor + t * (-2cm, -9cm);
            label(theorem, (0,3cm) - anchor);
            circle_NM ;
            jasper_render_segments ;
        \stopMPpage%
    }
    \startbuffer[BufA]
        \startframed[myframed]
            {\bf Theorem 2.2.2} A circle cannot have more than one 
            centre.\par {\it Proof.} Let us suggest that the conclusion 
            of the theorem is false; in particular that a circle has 
            two or more centres. Let us denote just two of these centres
            \im{O} and \im{P}.\par According to AXIOM 1, there is 
            exactly one line that passes through points \im{O} and 
            \im{P}. Let \im{M} and \im{N} be the points of intersection 
            of this line with the circle.\par Since we have suggested
            that \im{O} is a centre of the circle, the segments 
            \im{NO} and \im{OM} are congruent as two radii of the 
            circle centred at \im{O}: \im{NO=OM}. \color[white]{Similarly,
            \im{NP=PM} as two radii of the circle centred at 
            \im{P}.\par Point \im{O} lies in the interior of 
            \im{NP}, hence \im{NO<NP}; similarly, \im{PM<OM}.
            Then we can write: \dm{NO<NP=PM<OM=NO}, which means 
            that \im{NO<NO}. This result contradicts to the axiom 
            of congruence of segments which asserts that each 
            segments is congruent to itself. Thus the suggestion 
            of the existence of two (\im{P} and \im{O}) or more 
            distinct centres is false; therefore a circle cannot 
            have more than one centre, {\bf Q.E.D.}}
        \stopframed
    \stopbuffer
    % Since we have suggested
    % that \im{O} is a centre of the circle, the segments 
    % \im{NO} and \im{OM} are congruent as two radii of the 
    % circle centred at \im{O}: \im{NO=OM}.
    \dorecurse{11*30}{%
        \startMPpage
            render_theorem_box ;
            numeric t;
            t := 1;
            pair anchor;
            anchor := (1 - t) * center_anchor + t * topleft_anchor + t * (-2cm, -9cm);
            label(theorem, (0,3cm) - anchor);
            circle_NM ;
            jasper_render_segments ;
        \stopMPpage%
    }
    % line NO
    \dorecurse{3*30}{%
        \startMPpage
            render_theorem_box ;
            numeric t;
            t := 1;
            pair anchor;
            anchor := (1 - t) * center_anchor + t * topleft_anchor + t * (-2cm, -9cm);
            label(theorem, (0,3cm) - anchor);
            circle_NM ;
            jasper_curve[
                ustart="cx-R"
                ,ustop="cx"
                ,usamples="20"
                ,fx="u"
                ,fy="cy"
                ,fz="10"
                ,drawcommands="jasper_lightgraypenthick"
            ] ;
            jasper_point[
                fx="cx-R"
                ,fy="cy"
                ,fz="10"
                ,drawcommands="jasper_lightgray"
            ] ;
            jasper_point[
                fx="cx"
                ,fy="cy"
                ,fz="10"
                ,drawcommands="jasper_lightgray"
            ] ;
            jasper_render_segments ;
        \stopMPpage%
    }
    % shift to OM
    \dorecurse{1*30}{%
        \startMPpage
            render_theorem_box ;
            numeric t;
            t := 1;
            pair anchor;
            anchor := (1 - t) * center_anchor + t * topleft_anchor + t * (-2cm, -9cm);
            label(theorem, (0,3cm) - anchor);
            circle_NM ;
            jasper_curve[
                ustart="cx-R+R*#1/30"
                ,ustop="cx+R*#1/30"
                ,usamples="20"
                ,fx="u"
                ,fy="cy"
                ,fz="10"
                ,drawcommands="jasper_lightgraypenthick"
            ] ;
            jasper_point[
                fx="cx-R+R*#1/30"
                ,fy="cy"
                ,fz="10"
                ,drawcommands="jasper_lightgray"
            ] ;
            jasper_point[
                fx="cx+R*#1/30"
                ,fy="cy"
                ,fz="10"
                ,drawcommands="jasper_lightgray"
            ] ;
            jasper_render_segments ;
        \stopMPpage%
    }
    % line OM
    \dorecurse{4*30}{%
        \startMPpage
            render_theorem_box ;
            numeric t;
            t := 1;
            pair anchor;
            anchor := (1 - t) * center_anchor + t * topleft_anchor + t * (-2cm, -9cm);
            label(theorem, (0,3cm) - anchor);
            circle_NM ;
            jasper_curve[
                ustart="cx+R"
                ,ustop="cx"
                ,usamples="20"
                ,fx="u"
                ,fy="cy"
                ,fz="10"
                ,drawcommands="jasper_lightgraypenthick"
            ] ;
            jasper_point[
                fx="cx+R"
                ,fy="cy"
                ,fz="10"
                ,drawcommands="jasper_lightgray"
            ] ;
            jasper_point[
                fx="cx"
                ,fy="cy"
                ,fz="10"
                ,drawcommands="jasper_lightgray"
            ] ;
            jasper_render_segments ;
        \stopMPpage%
    }
    % pause with neither segment (second iteration)
    \dorecurse{3*30}{%
        \startMPpage
            render_theorem_box ;
            numeric t;
            t := 1;
            pair anchor;
            anchor := (1 - t) * center_anchor + t * topleft_anchor + t * (-2cm, -9cm);
            label(theorem, (0,3cm) - anchor);
            circle_NM ;
            jasper_render_segments ;
        \stopMPpage%
    }
    \startbuffer[BufA]
        \startframed[myframed]
            {\bf Theorem 2.2.2} A circle cannot have more than one 
            centre.\par {\it Proof.} Let us suggest that the conclusion 
            of the theorem is false; in particular that a circle has 
            two or more centres. Let us denote just two of these centres
            \im{O} and \im{P}.\par According to AXIOM 1, there is 
            exactly one line that passes through points \im{O} and 
            \im{P}. Let \im{M} and \im{N} be the points of intersection 
            of this line with the circle.\par Since we have suggested
            that \im{O} is a centre of the circle, the segments 
            \im{NO} and \im{OM} are congruent as two radii of the 
            circle centred at \im{O}: \im{NO=OM}. Similarly,
            \im{NP=PM} as two radii of the circle centred at 
            \im{P}.\par \color[white]{Point \im{O} lies in the interior of 
            \im{NP}, hence \im{NO<NP}; similarly, \im{PM<OM}.
            Then we can write: \dm{NO<NP=PM<OM=NO}, which means 
            that \im{NO<NO}. This result contradicts to the axiom 
            of congruence of segments which asserts that each 
            segments is congruent to itself. Thus the suggestion 
            of the existence of two (\im{P} and \im{O}) or more 
            distinct centres is false; therefore a circle cannot 
            have more than one centre, {\bf Q.E.D.}}
        \stopframed
    \stopbuffer
    % Similarly,
    % \im{NP=PM} as two radii of the circle centred at 
    % \im{P}.
    \dorecurse{8*30}{%
        \startMPpage
            render_theorem_box ;
            numeric t;
            t := 1;
            pair anchor;
            anchor := (1 - t) * center_anchor + t * topleft_anchor + t * (-2cm, -9cm);
            label(theorem, (0,3cm) - anchor);
            circle_NM ;
            jasper_render_segments ;
        \stopMPpage%
    }
    % line NP
    \dorecurse{2*30}{%
        \startMPpage
            render_theorem_box ;
            numeric t;
            t := 1;
            pair anchor;
            anchor := (1 - t) * center_anchor + t * topleft_anchor + t * (-2cm, -9cm);
            label(theorem, (0,3cm) - anchor);
            circle_NM ;
            jasper_curve[
                ustart="cx-R"
                ,ustop="cx+2"
                ,usamples="20"
                ,fx="u"
                ,fy="cy"
                ,fz="10"
                ,drawcommands="jasper_lightgraypenthick"
            ] ;
            jasper_point[
                fx="cx-R"
                ,fy="cy"
                ,fz="10"
                ,drawcommands="jasper_lightgray"
            ] ;
            jasper_point[
                fx="cx+2"
                ,fy="cy"
                ,fz="10"
                ,drawcommands="jasper_lightgray"
            ] ;
            jasper_render_segments ;
        \stopMPpage%
    }
    % interpolate
    \dorecurse{1*30}{%
        \startMPpage
            render_theorem_box ;
            numeric t;
            t := 1;
            pair anchor;
            anchor := (1 - t) * center_anchor + t * topleft_anchor + t * (-2cm, -9cm);
            label(theorem, (0,3cm) - anchor);
            circle_NM ;
            jasper_curve[
                ustart="cx-R+(R+2)*#1/30"
                ,ustop="cx+2+(R-2)*#1/30"
                ,usamples="20"
                ,fx="u"
                ,fy="cy"
                ,fz="10"
                ,drawcommands="jasper_lightgraypenthick"
            ] ;
            jasper_point[
                fx="cx-R+(R+2)*#1/30"
                ,fy="cy"
                ,fz="10"
                ,drawcommands="jasper_lightgray"
            ] ;
            jasper_point[
                fx="cx+2+(R-2)*#1/30"
                ,fy="cy"
                ,fz="10"
                ,drawcommands="jasper_lightgray"
            ] ;
            jasper_render_segments ;
        \stopMPpage%
    }
    % line PM
    \dorecurse{4*30}{%
        \startMPpage
            render_theorem_box ;
            numeric t;
            t := 1;
            pair anchor;
            anchor := (1 - t) * center_anchor + t * topleft_anchor + t * (-2cm, -9cm);
            label(theorem, (0,3cm) - anchor);
            circle_NM ;
            jasper_curve[
                ustart="cx+R"
                ,ustop="cx+2"
                ,usamples="20"
                ,fx="u"
                ,fy="cy"
                ,fz="10"
                ,drawcommands="jasper_lightgraypenthick"
            ] ;
            jasper_point[
                fx="cx+R"
                ,fy="cy"
                ,fz="10"
                ,drawcommands="jasper_lightgray"
            ] ;
            jasper_point[
                fx="cx+2"
                ,fy="cy"
                ,fz="10"
                ,drawcommands="jasper_lightgray"
            ] ;
            jasper_render_segments ;
        \stopMPpage%
    }
    % pause with neither segment (second iteration)
    \dorecurse{3*30}{%
        \startMPpage
            render_theorem_box ;
            numeric t;
            t := 1;
            pair anchor;
            anchor := (1 - t) * center_anchor + t * topleft_anchor + t * (-2cm, -9cm);
            label(theorem, (0,3cm) - anchor);
            circle_NM ;
            jasper_render_segments ;
        \stopMPpage%
    }
    \startbuffer[BufA]
        \startframed[myframed]
            {\bf Theorem 2.2.2} A circle cannot have more than one 
            centre.\par {\it Proof.} Let us suggest that the conclusion 
            of the theorem is false; in particular that a circle has 
            two or more centres. Let us denote just two of these centres
            \im{O} and \im{P}.\par According to AXIOM 1, there is 
            exactly one line that passes through points \im{O} and 
            \im{P}. Let \im{M} and \im{N} be the points of intersection 
            of this line with the circle.\par Since we have suggested
            that \im{O} is a centre of the circle, the segments 
            \im{NO} and \im{OM} are congruent as two radii of the 
            circle centred at \im{O}: \im{NO=OM}. Similarly,
            \im{NP=PM} as two radii of the circle centred at 
            \im{P}.\par Point \im{O} lies in the interior of 
            \im{NP}, hence \im{NO<NP}; \color[white]{similarly, \im{PM<OM}.
            Then we can write: \dm{NO<NP=PM<OM=NO}, which means 
            that \im{NO<NO}. This result contradicts to the axiom 
            of congruence of segments which asserts that each 
            segments is congruent to itself. Thus the suggestion 
            of the existence of two (\im{P} and \im{O}) or more 
            distinct centres is false; therefore a circle cannot 
            have more than one centre, {\bf Q.E.D.}}
        \stopframed
    \stopbuffer
    % Point \im{O} lies in the interior of 
    % \im{NP}, hence \im{NO<NP};
    \dorecurse{6*30}{%
        \startMPpage
            render_theorem_box ;
            numeric t;
            t := 1;
            pair anchor;
            anchor := (1 - t) * center_anchor + t * topleft_anchor + t * (-2cm, -9cm);
            label(theorem, (0,3cm) - anchor);
            circle_NM ;
            jasper_render_segments ;
        \stopMPpage%
    }
    % NP
    \dorecurse{2*30}{%
        \startMPpage
            render_theorem_box ;
            numeric t;
            t := 1;
            pair anchor;
            anchor := (1 - t) * center_anchor + t * topleft_anchor + t * (-2cm, -9cm);
            label(theorem, (0,3cm) - anchor);
            circle_NM ;
            jasper_curve[
                ustart="cx-R"
                ,ustop="cx+2"
                ,usamples="20"
                ,fx="u"
                ,fy="cy"
                ,fz="10"
                ,drawcommands="jasper_lightgraypenthick"
            ] ;
            jasper_point[
                fx="cx-R"
                ,fy="cy"
                ,fz="10"
                ,drawcommands="jasper_lightgray"
            ] ;
            jasper_point[
                fx="cx+2"
                ,fy="cy"
                ,fz="10"
                ,drawcommands="jasper_lightgray"
            ] ;
            jasper_render_segments ;
        \stopMPpage%
    }
    % NO<NP
    \dorecurse{1*30}{%
        \startMPpage
            render_theorem_box ;
            numeric t;
            t := 1;
            pair anchor;
            anchor := (1 - t) * center_anchor + t * topleft_anchor + t * (-2cm, -9cm);
            label(theorem, (0,3cm) - anchor);
            circle_NM ;
            jasper_curve[
                ustart="cx-R"
                ,ustop="cx+2-2*#1/30"
                ,usamples="20"
                ,fx="u"
                ,fy="cy"
                ,fz="10"
                ,drawcommands="jasper_lightgraypenthick"
            ] ;
            jasper_point[
                fx="cx-R"
                ,fy="cy"
                ,fz="10"
                ,drawcommands="jasper_lightgray"
            ] ;
            jasper_point[
                fx="cx+2-2*#1/30"
                ,fy="cy"
                ,fz="10"
                ,drawcommands="jasper_lightgray"
            ] ;
            jasper_render_segments ;
        \stopMPpage%
    }
     % NO
    \dorecurse{3*30}{%
        \startMPpage
            render_theorem_box ;
            numeric t;
            t := 1;
            pair anchor;
            anchor := (1 - t) * center_anchor + t * topleft_anchor + t * (-2cm, -9cm);
            label(theorem, (0,3cm) - anchor);
            circle_NM ;
            jasper_curve[
                ustart="cx-R"
                ,ustop="cx"
                ,usamples="20"
                ,fx="u"
                ,fy="cy"
                ,fz="10"
                ,drawcommands="jasper_lightgraypenthick"
            ] ;
            jasper_point[
                fx="cx-R"
                ,fy="cy"
                ,fz="10"
                ,drawcommands="jasper_lightgray"
            ] ;
            jasper_point[
                fx="cx"
                ,fy="cy"
                ,fz="10"
                ,drawcommands="jasper_lightgray"
            ] ;
            jasper_render_segments ;
        \stopMPpage%
    }
    % pause (no lines)
    \dorecurse{2*30}{%
        \startMPpage
            render_theorem_box ;
            numeric t;
            t := 1;
            pair anchor;
            anchor := (1 - t) * center_anchor + t * topleft_anchor + t * (-2cm, -9cm);
            label(theorem, (0,3cm) - anchor);
            circle_NM ;
            jasper_render_segments ;
        \stopMPpage%
    }
    \startbuffer[BufA]
        \startframed[myframed]
            {\bf Theorem 2.2.2} A circle cannot have more than one 
            centre.\par {\it Proof.} Let us suggest that the conclusion 
            of the theorem is false; in particular that a circle has 
            two or more centres. Let us denote just two of these centres
            \im{O} and \im{P}.\par According to AXIOM 1, there is 
            exactly one line that passes through points \im{O} and 
            \im{P}. Let \im{M} and \im{N} be the points of intersection 
            of this line with the circle.\par Since we have suggested
            that \im{O} is a centre of the circle, the segments 
            \im{NO} and \im{OM} are congruent as two radii of the 
            circle centred at \im{O}: \im{NO=OM}. Similarly,
            \im{NP=PM} as two radii of the circle centred at 
            \im{P}.\par Point \im{O} lies in the interior of 
            \im{NP}, hence \im{NO<NP}; similarly, \im{PM<OM}.
            \color[white]{Then we can write: \dm{NO<NP=PM<OM=NO}, which means 
            that \im{NO<NO}. This result contradicts to the axiom 
            of congruence of segments which asserts that each 
            segments is congruent to itself. Thus the suggestion 
            of the existence of two (\im{P} and \im{O}) or more 
            distinct centres is false; therefore a circle cannot 
            have more than one centre, {\bf Q.E.D.}}
        \stopframed
    \stopbuffer
    % similarly, \im{PM<OM}.
    \dorecurse{4*30}{%
        \startMPpage
            render_theorem_box ;
            numeric t;
            t := 1;
            pair anchor;
            anchor := (1 - t) * center_anchor + t * topleft_anchor + t * (-2cm, -9cm);
            label(theorem, (0,3cm) - anchor);
            circle_NM ;
            jasper_render_segments ;
        \stopMPpage%
    }
    % OM
    \dorecurse{2*30}{%
        \startMPpage
            render_theorem_box ;
            numeric t;
            t := 1;
            pair anchor;
            anchor := (1 - t) * center_anchor + t * topleft_anchor + t * (-2cm, -9cm);
            label(theorem, (0,3cm) - anchor);
            circle_NM ;
            jasper_curve[
                ustart="cx"
                ,ustop="cx+R"
                ,usamples="20"
                ,fx="u"
                ,fy="cy"
                ,fz="10"
                ,drawcommands="jasper_lightgraypenthick"
            ] ;
            jasper_point[
                fx="cx"
                ,fy="cy"
                ,fz="10"
                ,drawcommands="jasper_lightgray"
            ] ;
            jasper_point[
                fx="cx+R"
                ,fy="cy"
                ,fz="10"
                ,drawcommands="jasper_lightgray"
            ] ;
            jasper_render_segments ;
        \stopMPpage%
    }
    % PM<OM
    \dorecurse{1*30}{%
        \startMPpage
            render_theorem_box ;
            numeric t;
            t := 1;
            pair anchor;
            anchor := (1 - t) * center_anchor + t * topleft_anchor + t * (-2cm, -9cm);
            label(theorem, (0,3cm) - anchor);
            circle_NM ;
            jasper_curve[
                ustart="cx+2*#1/30"
                ,ustop="cx+R"
                ,usamples="20"
                ,fx="u"
                ,fy="cy"
                ,fz="10"
                ,drawcommands="jasper_lightgraypenthick"
            ] ;
            jasper_point[
                fx="cx+2*#1/30"
                ,fy="cy"
                ,fz="10"
                ,drawcommands="jasper_lightgray"
            ] ;
            jasper_point[
                fx="cx+R"
                ,fy="cy"
                ,fz="10"
                ,drawcommands="jasper_lightgray"
            ] ;
            jasper_render_segments ;
        \stopMPpage%
    }
    % PM
    \dorecurse{3*30}{%
        \startMPpage
            render_theorem_box ;
            numeric t;
            t := 1;
            pair anchor;
            anchor := (1 - t) * center_anchor + t * topleft_anchor + t * (-2cm, -9cm);
            label(theorem, (0,3cm) - anchor);
            circle_NM ;
            jasper_curve[
                ustart="cx+2"
                ,ustop="cx+R"
                ,usamples="20"
                ,fx="u"
                ,fy="cy"
                ,fz="10"
                ,drawcommands="jasper_lightgraypenthick"
            ] ;
            jasper_point[
                fx="cx+2"
                ,fy="cy"
                ,fz="10"
                ,drawcommands="jasper_lightgray"
            ] ;
            jasper_point[
                fx="cx+R"
                ,fy="cy"
                ,fz="10"
                ,drawcommands="jasper_lightgray"
            ] ;
            jasper_render_segments ;
        \stopMPpage%
    }
    % pause
    \dorecurse{3*30}{%
        \startMPpage
            render_theorem_box ;
            numeric t;
            t := 1;
            pair anchor;
            anchor := (1 - t) * center_anchor + t * topleft_anchor + t * (-2cm, -9cm);
            label(theorem, (0,3cm) - anchor);
            circle_NM ;
            jasper_render_segments ;
        \stopMPpage%
    }
    \startbuffer[BufA]
        \startframed[myframed]
            {\bf Theorem 2.2.2} A circle cannot have more than one 
            centre.\par {\it Proof.} Let us suggest that the conclusion 
            of the theorem is false; in particular that a circle has 
            two or more centres. Let us denote just two of these centres
            \im{O} and \im{P}.\par According to AXIOM 1, there is 
            exactly one line that passes through points \im{O} and 
            \im{P}. Let \im{M} and \im{N} be the points of intersection 
            of this line with the circle.\par Since we have suggested
            that \im{O} is a centre of the circle, the segments 
            \im{NO} and \im{OM} are congruent as two radii of the 
            circle centred at \im{O}: \im{NO=OM}. Similarly,
            \im{NP=PM} as two radii of the circle centred at 
            \im{P}.\par Point \im{O} lies in the interior of 
            \im{NP}, hence \im{NO<NP}; similarly, \im{PM<OM}.
            Then we can write: \dm{NO<NP=PM<OM=NO}, \color[white]{which means 
            that \im{NO<NO}. This result contradicts to the axiom 
            of congruence of segments which asserts that each 
            segments is congruent to itself. Thus the suggestion 
            of the existence of two (\im{P} and \im{O}) or more 
            distinct centres is false; therefore a circle cannot 
            have more than one centre, {\bf Q.E.D.}}
        \stopframed
    \stopbuffer
    % Then we can write: \dm{NO<NP=PM<OM=NO},
    \dorecurse{10*30}{%
        \startMPpage
            render_theorem_box ;
            numeric t;
            t := 1;
            pair anchor;
            anchor := (1 - t) * center_anchor + t * topleft_anchor + t * (-2cm, -9cm);
            label(theorem, (0,3cm) - anchor);
            circle_NM ;
            jasper_render_segments ;
        \stopMPpage%
    }
    % NP
    \dorecurse{2*30}{%
        \startMPpage
            render_theorem_box ;
            numeric t;
            t := 1;
            pair anchor;
            anchor := (1 - t) * center_anchor + t * topleft_anchor + t * (-2cm, -9cm);
            label(theorem, (0,3cm) - anchor);
            circle_NM ;
            jasper_curve[
                ustart="cx-R"
                ,ustop="cx+2"
                ,usamples="20"
                ,fx="u"
                ,fy="cy"
                ,fz="10"
                ,drawcommands="jasper_lightgraypenthick"
            ] ;
            jasper_point[
                fx="cx-R"
                ,fy="cy"
                ,fz="10"
                ,drawcommands="jasper_lightgray"
            ] ;
            jasper_point[
                fx="cx+2"
                ,fy="cy"
                ,fz="10"
                ,drawcommands="jasper_lightgray"
            ] ;
            jasper_render_segments ;
        \stopMPpage%
    }
    % NO<NP
    \dorecurse{1*30}{%
        \startMPpage
            render_theorem_box ;
            numeric t;
            t := 1;
            pair anchor;
            anchor := (1 - t) * center_anchor + t * topleft_anchor + t * (-2cm, -9cm);
            label(theorem, (0,3cm) - anchor);
            circle_NM ;
            jasper_curve[
                ustart="cx-R"
                ,ustop="cx+2-2*#1/30"
                ,usamples="20"
                ,fx="u"
                ,fy="cy"
                ,fz="10"
                ,drawcommands="jasper_lightgraypenthick"
            ] ;
            jasper_point[
                fx="cx-R"
                ,fy="cy"
                ,fz="10"
                ,drawcommands="jasper_lightgray"
            ] ;
            jasper_point[
                fx="cx+2-2*#1/30"
                ,fy="cy"
                ,fz="10"
                ,drawcommands="jasper_lightgray"
            ] ;
            jasper_render_segments ;
        \stopMPpage%
    }
    % NO
    \dorecurse{2*30}{%
        \startMPpage
            render_theorem_box ;
            numeric t;
            t := 1;
            pair anchor;
            anchor := (1 - t) * center_anchor + t * topleft_anchor + t * (-2cm, -9cm);
            label(theorem, (0,3cm) - anchor);
            circle_NM ;
            jasper_curve[
                ustart="cx-R"
                ,ustop="cx"
                ,usamples="20"
                ,fx="u"
                ,fy="cy"
                ,fz="10"
                ,drawcommands="jasper_lightgraypenthick"
            ] ;
            jasper_point[
                fx="cx-R"
                ,fy="cy"
                ,fz="10"
                ,drawcommands="jasper_lightgray"
            ] ;
            jasper_point[
                fx="cx"
                ,fy="cy"
                ,fz="10"
                ,drawcommands="jasper_lightgray"
            ] ;
            jasper_render_segments ;
        \stopMPpage%
    }
    % NO=OM
    \dorecurse{1*30}{%
        \startMPpage
            render_theorem_box ;
            numeric t;
            t := 1;
            pair anchor;
            anchor := (1 - t) * center_anchor + t * topleft_anchor + t * (-2cm, -9cm);
            label(theorem, (0,3cm) - anchor);
            circle_NM ;
            jasper_curve[
                ustart="cx-R+7*#1/30"
                ,ustop="cx+7*#1/30"
                ,usamples="20"
                ,fx="u"
                ,fy="cy"
                ,fz="10"
                ,drawcommands="jasper_lightgraypenthick"
            ] ;
            jasper_point[
                fx="cx-R+7*#1/30"
                ,fy="cy"
                ,fz="10"
                ,drawcommands="jasper_lightgray"
            ] ;
            jasper_point[
                fx="cx+7*#1/30"
                ,fy="cy"
                ,fz="10"
                ,drawcommands="jasper_lightgray"
            ] ;
            jasper_render_segments ;
        \stopMPpage%
    }
    % OM
    \dorecurse{2*30}{%
        \startMPpage
            render_theorem_box ;
            numeric t;
            t := 1;
            pair anchor;
            anchor := (1 - t) * center_anchor + t * topleft_anchor + t * (-2cm, -9cm);
            label(theorem, (0,3cm) - anchor);
            circle_NM ;
            jasper_curve[
                ustart="cx"
                ,ustop="cx+R"
                ,usamples="20"
                ,fx="u"
                ,fy="cy"
                ,fz="10"
                ,drawcommands="jasper_lightgraypenthick"
            ] ;
            jasper_point[
                fx="cx"
                ,fy="cy"
                ,fz="10"
                ,drawcommands="jasper_lightgray"
            ] ;
            jasper_point[
                fx="cx+R"
                ,fy="cy"
                ,fz="10"
                ,drawcommands="jasper_lightgray"
            ] ;
            jasper_render_segments ;
        \stopMPpage%
    }
    % PM<OM
    \dorecurse{1*30}{%
        \startMPpage
            render_theorem_box ;
            numeric t;
            t := 1;
            pair anchor;
            anchor := (1 - t) * center_anchor + t * topleft_anchor + t * (-2cm, -9cm);
            label(theorem, (0,3cm) - anchor);
            circle_NM ;
            jasper_curve[
                ustart="cx+2*#1/30"
                ,ustop="cx+R"
                ,usamples="20"
                ,fx="u"
                ,fy="cy"
                ,fz="10"
                ,drawcommands="jasper_lightgraypenthick"
            ] ;
            jasper_point[
                fx="cx+2*#1/30"
                ,fy="cy"
                ,fz="10"
                ,drawcommands="jasper_lightgray"
            ] ;
            jasper_point[
                fx="cx+R"
                ,fy="cy"
                ,fz="10"
                ,drawcommands="jasper_lightgray"
            ] ;
            jasper_render_segments ;
        \stopMPpage%
    }
    % PM
    \dorecurse{2*30}{%
        \startMPpage
            render_theorem_box ;
            numeric t;
            t := 1;
            pair anchor;
            anchor := (1 - t) * center_anchor + t * topleft_anchor + t * (-2cm, -9cm);
            label(theorem, (0,3cm) - anchor);
            circle_NM ;
            jasper_curve[
                ustart="cx+2"
                ,ustop="cx+R"
                ,usamples="20"
                ,fx="u"
                ,fy="cy"
                ,fz="10"
                ,drawcommands="jasper_lightgraypenthick"
            ] ;
            jasper_point[
                fx="cx+2"
                ,fy="cy"
                ,fz="10"
                ,drawcommands="jasper_lightgray"
            ] ;
            jasper_point[
                fx="cx+R"
                ,fy="cy"
                ,fz="10"
                ,drawcommands="jasper_lightgray"
            ] ;
            jasper_render_segments ;
        \stopMPpage%
    }%
    % PM<NO and PM=NO???
    \dorecurse{1*30}{%
        \startMPpage
            render_theorem_box ;
            numeric t;
            t := 1;
            pair anchor;
            anchor := (1 - t) * center_anchor + t * topleft_anchor + t * (-2cm, -9cm);
            label(theorem, (0,3cm) - anchor);
            circle_NM ;
            jasper_curve[
                ustart="cx+2-(R+2)*#1/30"
                ,ustop="cx+R-(R+2)*#1/30"
                ,usamples="20"
                ,fx="u"
                ,fy="cy"
                ,fz="10"
                ,drawcommands="jasper_lightgraypenthick"
            ] ;
            jasper_point[
                fx="cx+2-(R+2)*#1/30"
                ,fy="cy"
                ,fz="10"
                ,drawcommands="jasper_lightgray"
            ] ;
            jasper_point[
                fx="cx+R-(R+2)*#1/30"
                ,fy="cy"
                ,fz="10"
                ,drawcommands="jasper_lightgray"
            ] ;
            jasper_render_segments ;
        \stopMPpage%
    }%
    \startMPinclusions
        def labeled_circle =
            circle_NM ;
            jasper_curve[
                ustart="cx-R"
                ,ustop="cx-2"
                ,usamples="20"
                ,fx="u"
                ,fy="cy"
                ,fz="10"
                ,drawcommands="jasper_lightgraypenthick"
            ] ;
            jasper_point[
                fx="cx-R"
                ,fy="cy"
                ,fz="10"
                ,drawcommands="jasper_lightgray"
            ] ;
            jasper_point[
                fx="cx-2"
                ,fy="cy"
                ,fz="10"
                ,drawcommands="jasper_lightgray"
            ] ;
        enddef ;
    \stopMPinclusions%
    % -> NO<NO
    \dorecurse{3*30}{%
        \startMPpage
            render_theorem_box ;
            numeric t;
            t := 1;
            pair anchor;
            anchor := (1 - t) * center_anchor + t * topleft_anchor + t * (-2cm, -9cm);
            label(theorem, (0,3cm) - anchor);
            labeled_circle ;
            jasper_render_segments ;
        \stopMPpage%
    }
    \startbuffer[BufA]
        \startframed[myframed]
            {\bf Theorem 2.2.2} A circle cannot have more than one 
            centre.\par {\it Proof.} Let us suggest that the conclusion 
            of the theorem is false; in particular that a circle has 
            two or more centres. Let us denote just two of these centres
            \im{O} and \im{P}.\par According to AXIOM 1, there is 
            exactly one line that passes through points \im{O} and 
            \im{P}. Let \im{M} and \im{N} be the points of intersection 
            of this line with the circle.\par Since we have suggested
            that \im{O} is a centre of the circle, the segments 
            \im{NO} and \im{OM} are congruent as two radii of the 
            circle centred at \im{O}: \im{NO=OM}. Similarly,
            \im{NP=PM} as two radii of the circle centred at 
            \im{P}.\par Point \im{O} lies in the interior of 
            \im{NP}, hence \im{NO<NP}; similarly, \im{PM<OM}.
            Then we can write: \dm{NO<NP=PM<OM=NO}, which means 
            that \im{NO<NO}. \color[white]{This result contradicts to the axiom 
            of congruence of segments which asserts that each 
            segments is congruent to itself. Thus the suggestion 
            of the existence of two (\im{P} and \im{O}) or more 
            distinct centres is false; therefore a circle cannot 
            have more than one centre, {\bf Q.E.D.}}
        \stopframed
    \stopbuffer
    % which means 
    % that \im{NO<NO}.
    \dorecurse{5*30}{%
        \startMPpage
            render_theorem_box ;
            numeric t;
            t := 1;
            pair anchor;
            anchor := (1 - t) * center_anchor + t * topleft_anchor + t * (-2cm, -9cm);
            label(theorem, (0,3cm) - anchor);
            labeled_circle ;
            jasper_render_segments ;
        \stopMPpage%
    }
    \startbuffer[BufA]
        \startframed[myframed]
            {\bf Theorem 2.2.2} A circle cannot have more than one 
            centre.\par {\it Proof.} Let us suggest that the conclusion 
            of the theorem is false; in particular that a circle has 
            two or more centres. Let us denote just two of these centres
            \im{O} and \im{P}.\par According to AXIOM 1, there is 
            exactly one line that passes through points \im{O} and 
            \im{P}. Let \im{M} and \im{N} be the points of intersection 
            of this line with the circle.\par Since we have suggested
            that \im{O} is a centre of the circle, the segments 
            \im{NO} and \im{OM} are congruent as two radii of the 
            circle centred at \im{O}: \im{NO=OM}. Similarly,
            \im{NP=PM} as two radii of the circle centred at 
            \im{P}.\par Point \im{O} lies in the interior of 
            \im{NP}, hence \im{NO<NP}; similarly, \im{PM<OM}.
            Then we can write: \dm{NO<NP=PM<OM=NO}, which means 
            that \im{NO<NO}. This result contradicts to the axiom 
            of congruence of segments which asserts that each 
            segments is congruent to itself. \color[white]{Thus the suggestion 
            of the existence of two (\im{P} and \im{O}) or more 
            distinct centres is false; therefore a circle cannot 
            have more than one centre, {\bf Q.E.D.}}
        \stopframed
    \stopbuffer
    % This result contradicts to the axiom 
    % of congruence of segments which asserts that each 
    % segments is congruent to itself.
    \dorecurse{10*30}{%
        \startMPpage
            render_theorem_box ;
            numeric t;
            t := 1;
            pair anchor;
            anchor := (1 - t) * center_anchor + t * topleft_anchor + t * (-2cm, -9cm);
            label(theorem, (0,3cm) - anchor);
            labeled_circle ;
            jasper_render_segments ;
        \stopMPpage%
    }
    \startbuffer[BufA]
        \startframed[myframed]
            {\bf Theorem 2.2.2} A circle cannot have more than one 
            centre.\par {\it Proof.} Let us suggest that the conclusion 
            of the theorem is false; in particular that a circle has 
            two or more centres. Let us denote just two of these centres
            \im{O} and \im{P}.\par According to AXIOM 1, there is 
            exactly one line that passes through points \im{O} and 
            \im{P}. Let \im{M} and \im{N} be the points of intersection 
            of this line with the circle.\par Since we have suggested
            that \im{O} is a centre of the circle, the segments 
            \im{NO} and \im{OM} are congruent as two radii of the 
            circle centred at \im{O}: \im{NO=OM}. Similarly,
            \im{NP=PM} as two radii of the circle centred at 
            \im{P}.\par Point \im{O} lies in the interior of 
            \im{NP}, hence \im{NO<NP}; similarly, \im{PM<OM}.
            Then we can write: \dm{NO<NP=PM<OM=NO}, which means 
            that \im{NO<NO}. This result contradicts to the axiom 
            of congruence of segments which asserts that each 
            segments is congruent to itself. Thus the suggestion 
            of the existence of two (\im{P} and \im{O}) or more 
            distinct centres is false; \color[white]{therefore a circle cannot 
            have more than one centre, {\bf Q.E.D.}}
        \stopframed
    \stopbuffer
    % Thus the suggestion 
    % of the existence of two (\im{P} and \im{O}) or more 
    % distinct centres is false; 
    \dorecurse{9*30}{%
        \startMPpage
            render_theorem_box ;
            numeric t;
            t := 1;
            pair anchor;
            anchor := (1 - t) * center_anchor + t * topleft_anchor + t * (-2cm, -9cm);
            label(theorem, (0,3cm) - anchor);
            labeled_circle ;
            jasper_render_segments ;
        \stopMPpage%
    }
    \startbuffer[BufA]
        \startframed[myframed]
            {\bf Theorem 2.2.2} A circle cannot have more than one 
            centre.\par {\it Proof.} Let us suggest that the conclusion 
            of the theorem is false; in particular that a circle has 
            two or more centres. Let us denote just two of these centres
            \im{O} and \im{P}.\par According to AXIOM 1, there is 
            exactly one line that passes through points \im{O} and 
            \im{P}. Let \im{M} and \im{N} be the points of intersection 
            of this line with the circle.\par Since we have suggested
            that \im{O} is a centre of the circle, the segments 
            \im{NO} and \im{OM} are congruent as two radii of the 
            circle centred at \im{O}: \im{NO=OM}. Similarly,
            \im{NP=PM} as two radii of the circle centred at 
            \im{P}.\par Point \im{O} lies in the interior of 
            \im{NP}, hence \im{NO<NP}; similarly, \im{PM<OM}.
            Then we can write: \dm{NO<NP=PM<OM=NO}, which means 
            that \im{NO<NO}. This result contradicts to the axiom 
            of congruence of segments which asserts that each 
            segments is congruent to itself. Thus the suggestion 
            of the existence of two (\im{P} and \im{O}) or more 
            distinct centres is false; therefore a circle cannot 
            have more than one centre, \color[white]{{\bf Q.E.D.}}
        \stopframed
    \stopbuffer
    % therefore a circle cannot 
    % have more than one centre,
    \dorecurse{5*30}{%
        \startMPpage
            render_theorem_box ;
            numeric t;
            t := 1;
            pair anchor;
            anchor := (1 - t) * center_anchor + t * topleft_anchor + t * (-2cm, -9cm);
            label(theorem, (0,3cm) - anchor);
            labeled_circle ;
            jasper_render_segments ;
        \stopMPpage%
    }
    \startbuffer[BufA]
        \startframed[myframed]
            {\bf Theorem 2.2.2} A circle cannot have more than one 
            centre.\par {\it Proof.} Let us suggest that the conclusion 
            of the theorem is false; in particular that a circle has 
            two or more centres. Let us denote just two of these centres
            \im{O} and \im{P}.\par According to AXIOM 1, there is 
            exactly one line that passes through points \im{O} and 
            \im{P}. Let \im{M} and \im{N} be the points of intersection 
            of this line with the circle.\par Since we have suggested
            that \im{O} is a centre of the circle, the segments 
            \im{NO} and \im{OM} are congruent as two radii of the 
            circle centred at \im{O}: \im{NO=OM}. Similarly,
            \im{NP=PM} as two radii of the circle centred at 
            \im{P}.\par Point \im{O} lies in the interior of 
            \im{NP}, hence \im{NO<NP}; similarly, \im{PM<OM}.
            Then we can write: \dm{NO<NP=PM<OM=NO}, which means 
            that \im{NO<NO}. This result contradicts to the axiom 
            of congruence of segments which asserts that each 
            segments is congruent to itself. Thus the suggestion 
            of the existence of two (\im{P} and \im{O}) or more 
            distinct centres is false; therefore a circle cannot 
            have more than one centre, {\bf Q.E.D.}
        \stopframed
    \stopbuffer
    % {\bf Q.E.D.}
    \dorecurse{15*30}{%
        \startMPpage
            render_theorem_box ;
            numeric t;
            t := 1;
            pair anchor;
            anchor := (1 - t) * center_anchor + t * topleft_anchor + t * (-2cm, -9cm);
            label(theorem, (0,3cm) - anchor);
            labeled_circle ;
            jasper_render_segments ;
        \stopMPpage%
    }
\stoptext
